\lecture{41}{Децентрализованный SGD: скорость и устойчивость}

\sourcevideo
    {https://www.youtube.com/playlist?list=PLOzRYVm0a65fhKDS8cYIC7YCuVGJ4FOJx}
    {Week 11: Lecture 41: Decentralized Stochastic Gradient Descent -2}

В этой лекции изучается влияние коммуникационного графа на скорость
децентрализованного стохастического градиентного спуска. Отдельно
рассматривается переходный период до появления линейного ускорения и
способ ограничить влияние агента, отклоняющегося от протокола.

\section{Модель и предположения}

Пусть \(N\) агентов решают задачу
\[
    \min_{x\in\mathbb R^d}F(x),
    \qquad
    F(x)=\frac1N\sum_{i=1}^{N}f_i(x),
\]
где
\[
    f_i(x)
    =
    \mathbb E_{\xi\sim\mathcal D_i}
    \bigl[\ell(x;\xi)\bigr].
\]
Агент \(i\) хранит локальную модель \(x_i^k\) и вычисляет
стохастический градиент
\[
    g_i^k=g_i(x_i^k;\xi_i^k).
\]
Один шаг децентрализованного SGD имеет вид
\[
    x_i^{k+\frac12}
    =
    x_i^k-\gamma g_i^k,
    \qquad
    x_i^{k+1}
    =
    \sum_{j=1}^{N}w_{ij}x_j^{k+\frac12}.
\]
Матрица смешивания \(W=(w_{ij})\) согласована с графом и предполагается
дважды стохастической.

Каждая \(f_i\) имеет \(L\)-липшицев градиент:
\[
    \lVert\nabla f_i(x)-\nabla f_i(y)\rVert
    \le
    L\lVert x-y\rVert.
\]
Стохастические градиенты несмещены,
\[
    \mathbb E[g_i^k\mid x_i^k]
    =
    \nabla f_i(x_i^k),
\]
и имеют ограниченную условную дисперсию:
\[
    \mathbb E
    \left[
        \lVert
            g_i^k-\nabla f_i(x_i^k)
        \rVert^2
        \,\middle|\,
        x_i^k
    \right]
    \le
    \sigma^2.
\]
Шумы различных агентов считаются независимыми.

Неоднородность локальных данных измеряется условием
\[
    \frac1N
    \sum_{i=1}^{N}
    \lVert
        \nabla f_i(x)-\nabla F(x)
    \rVert^2
    \le
    b^2.
\]
При \(b=0\) локальные градиенты совпадают с глобальным. Для конечных
эмпирических выборок это является идеализированным случаем даже тогда,
когда данные получены из одного распределения.

\section{Смешивание и оценка скорости}

Положим
\[
    J=\frac1N\mathbf1\mathbf1^{\mathsf T},
    \qquad
    \rho=\lVert W-J\rVert_2.
\]
Для связной апериодической сети
\[
    0\le\rho<1.
\]
Малое \(\rho\) соответствует быстрому смешиванию, а величина
\(1-\rho\) играет роль спектрального разрыва.

Обозначим среднюю модель через
\[
    \bar x^k
    =
    \frac1N\sum_{i=1}^{N}x_i^k.
\]
Для шага порядка \(\gamma\asymp T^{-1/2}\) получаем оценку
следующей структуры:
\[
\begin{aligned}
    \frac1T
    \sum_{k=0}^{T-1}
    \mathbb E
    \lVert\nabla F(\bar x^k)\rVert^2
    \lesssim{}&
    \frac{\sigma}{\sqrt{NT}}
    \\
    &+
    \frac{
        \rho^{2/3}\sigma^{2/3}
    }{
        (1-\rho)^{1/3}T^{2/3}
    }
    \\
    &+
    \frac{
        \rho^{2/3}b^{2/3}
    }{
        (1-\rho)^{2/3}T^{2/3}
    }.
\end{aligned}
\]
Знак \(\lesssim\) скрывает численные константы, гладкость и начальный
функциональный разрыв.

Первое слагаемое совпадает с ведущим стохастическим членом
параллельного SGD. Второе возникает из-за отсутствия точной глобальной
синхронизации, третье — из-за различия локальных функций. Сетевые
члены убывают как \(T^{-2/3}\), то есть асимптотически быстрее, чем
\(T^{-1/2}\). Поэтому при достаточно большом числе итераций ведущим
становится
\[
    O\left(\frac{\sigma}{\sqrt{NT}}\right),
\]
и возникает линейное ускорение по числу агентов.

Для дальнейшего анализа используются порядки \(T^{-1/2}\), \(T^{-2/3}\) и зависимости от \(\rho\), \(1-\rho\), \(\sigma\) и \(b\); точные константы не влияют на качественное сравнение режимов.

\section{Переходный период}

Переходный период определяется моментом, когда сетевые члены становятся
не больше ведущего стохастического слагаемого. При однородных данных,
то есть \(b=0\), сравнение первых двух членов даёт порядок
\[
    T_{\mathrm{tr}}
    =
    O\left(
        \frac{\rho^4N^3}{(1-\rho)^2}
    \right),
\]
если опустить зависимость от \(\sigma\) и других констант.

При неоднородных данных порядок ухудшается до
\[
    T_{\mathrm{tr}}
    =
    O\left(
        \frac{\rho^4N^3}{(1-\rho)^4}
    \right).
\]
Следовательно, линейное ускорение является асимптотическим свойством:
на начальном участке сетевое рассогласование может доминировать.
Переходный период особенно велик при \(\rho\), близком к единице.

Коммуникационная стоимость одной итерации зависит от максимальной
степени графа. Если она равна \(d_{\max}\), то один агент передаёт
порядка
\[
    O(d_{\max}d)
\]
чисел. Поэтому требуется одновременно ограничивать степень графа и
обеспечивать малое значение \(\rho\).

Полный граф быстро смешивает информацию, но имеет степень порядка
\(N\). Кольцо имеет постоянную степень, однако
\[
    1-\rho=\Theta(N^{-2}),
\]
из-за чего его переходный период быстро растёт.

\section{Статический экспоненциальный граф}

Пронумеруем вершины числами
\[
    0,\ldots,N-1.
\]
Из вершины \(i\) проводятся дуги к вершинам
\[
    i+2^r\pmod N,
    \qquad
    r=0,\ldots,\lceil\log_2N\rceil-1.
\]
Таким образом, агент соединяется с вершинами на циклических
расстояниях
\[
    1,2,4,8,\ldots.
\]
Исходящая степень имеет порядок
\[
    O(\log N).
\]

Пусть используется \(q\) ненулевых смещений. Один из вариантов
матрицы смешивания задаётся формулой
\[
    w_{ij}
    =
    \begin{cases}
        \dfrac1{q+1},
        &
        j=i,
        \\[2mm]
        \dfrac1{q+1},
        &
        j=i+2^r\pmod N
        \text{ для некоторого }r,
        \\[2mm]
        0,
        &
        \text{иначе}.
    \end{cases}
\]
Благодаря циклической симметрии матрица является строчно- и
столбцово-стохастической.

Экспоненциальный граф распространяет информацию существенно быстрее
пути или кольца, сохраняя степень порядка \(\log N\). В лекции он
выделяется как удачный компромисс между стоимостью одной итерации и
переходным периодом, но оптимальность среди всех возможных топологий
не утверждается.

\section{Отклоняющийся агент}

Рассмотрим локальные функции
\[
    f_i(x)=\frac12(x-i)^2,
    \qquad
    i=1,\ldots,5.
\]
Глобальный минимум равен
\[
    x^\star=3.
\]
Пусть агенты \(1,\ldots,4\) выполняют протокол, а агент \(5\)
игнорирует согласование и минимизирует только
\[
    f_5(x)=\frac12(x-5)^2.
\]
Его обновление имеет вид
\[
    x_5^{k+1}
    =
    x_5^k-\eta(x_5^k-5),
\]
поэтому он стремится к точке \(5\), продолжая передавать соседям своё
состояние. Обычное усреднение способно сместить честных агентов от
социального решения \(3\).

Это более узкая модель, чем полный византийский противник: агент
отклоняется от протокола, но не посылает различным соседям
произвольные и разные сообщения.

\section{Отсечённое смешивание}

Для \(z\in\mathbb R^d\) и \(\tau>0\) определим
\[
    \operatorname{clip}(z,\tau)
    =
    \min
    \left\{
        1,
        \frac{\tau}{\lVert z\rVert}
    \right\}z
\]
при \(z\ne0\), а
\[
    \operatorname{clip}(0,\tau)=0.
\]
Тогда
\[
    \lVert\operatorname{clip}(z,\tau)\rVert
    \le\tau.
\]
При \(\lVert z\rVert\le\tau\) вектор не изменяется; при большей норме
сохраняется его направление, но ограничивается величина.

Вместо обычного gossip агент вычисляет
\[
    z_i^k
    =
    x_i^k
    +
    \sum_{j\in\mathcal N_i}
    w_{ij}
    \operatorname{clip}
    \left(
        x_j^k-x_i^k,
        \tau
    \right)
\]
и затем выполняет локальный шаг
\[
    x_i^{k+1}
    =
    z_i^k-\eta\nabla f_i(z_i^k).
\]
Вклад одного соседа в смешанное состояние ограничен величиной,
пропорциональной \(\tau\), поэтому удалённая локальная модель не может
иметь неограниченного влияния.

При большом \(\tau\) алгоритм близок к обычному усреднению и слабо
ограничивает отклоняющегося агента. Малое \(\tau\) повышает
устойчивость, но замедляет согласование честных агентов и может внести
смещение. Численный пример лекции показывает улучшение поведения
честных агентов при уменьшении порога, но сам по себе не является
общей теоремой.

Для строгой защиты от византийских противников необходимы
дополнительные предположения о числе противников, связности графа,
локальных функциях, пороге отсечения и модели передаваемых сообщений.
Одно только отсечение таких гарантий не даёт.
